\documentclass[11pt]{amsart}

% Encoding and fonts
\usepackage[T1]{fontenc}
\usepackage[utf8]{inputenc}
\usepackage{lmodern}

% Math packages
\usepackage{amsmath}
\usepackage{amssymb}
\usepackage{amsfonts}
\usepackage{amsthm}
\usepackage{mathtools}

% Layout and utilities
\usepackage{geometry}
\usepackage{enumitem}
\usepackage[hidelinks]{hyperref}

\geometry{margin=1in}

\numberwithin{equation}{section}
\setcounter{tocdepth}{3}
\setcounter{secnumdepth}{3}

% ---------------------------------------------------------
% Theorem-like environments
% ---------------------------------------------------------

\theoremstyle{plain}
\newtheorem{theorem}{Theorem}[section]
\newtheorem{proposition}[theorem]{Proposition}
\newtheorem{lemma}[theorem]{Lemma}
\newtheorem{corollary}[theorem]{Corollary}
\newtheorem{conjecture}[theorem]{Conjecture}

\theoremstyle{definition}
\newtheorem{definition}[theorem]{Definition}
\newtheorem{experiment}[theorem]{Experiment}
\newtheorem{heuristic}[theorem]{Heuristic}
\newtheorem{question}[theorem]{Question}

\theoremstyle{remark}
\newtheorem{remark}[theorem]{Remark}
\newtheorem{warning}[theorem]{Warning}

\title{Exponential sums
of Zeta function zeros}

\author{}
\date{}

\begin{document}

\maketitle

\tableofcontents

\bigskip


\section{The Riemann Zeta Function}
\label{sec:zeta_function}

This section contains standard material.
The classical references are
\cite{Titchmarsh1986},
\cite{Edwards2001},
and
\cite{IwaniecKowalski2004}.

\subsection{Definition}
\label{subsec:zeta_definition}

\begin{definition}
\label{def:zeta_function}

The Riemann zeta function is defined for
$\Re(s)>1$
by the absolutely convergent series

\begin{equation}
\zeta(s)
=
\sum_{n=1}^{\infty}
\frac{1}{n^{s}}.
\label{eq:zeta_dirichlet_series}
\end{equation}

\end{definition}

Euler showed that the same function admits the product expansion

\begin{equation}
\zeta(s)
=
\prod_{p}
\left(
1-p^{-s}
\right)^{-1},
\qquad
\Re(s)>1,
\label{eq:euler_product}
\end{equation}

where the product extends over all prime numbers.
Equation
\eqref{eq:euler_product}
establishes the fundamental connection between
$\zeta(s)$
and the distribution of primes.

\subsection{Analytic Continuation}
\label{subsec:analytic_continuation}

A classical theorem of Riemann states that
$\zeta(s)$
admits a meromorphic continuation to the whole complex plane,
with a unique simple pole at

\begin{equation}
s=1.
\label{eq:zeta_pole}
\end{equation}

See
\cite{Edwards2001}
or
\cite{Titchmarsh1986}.

\subsection{Functional Equation}
\label{subsec:functional_equation}

Define

\begin{equation}
\xi(s)
=
\frac12
s(s-1)
\pi^{-s/2}
\Gamma\!\left(\frac{s}{2}\right)
\zeta(s).
\label{eq:xi_definition}
\end{equation}

The completed zeta function satisfies

\begin{equation}
\xi(s)
=
\xi(1-s).
\label{eq:functional_equation}
\end{equation}

Equivalently,

\begin{equation}
\zeta(s)
=
2^{s}
\pi^{\,s-1}
\sin\!\left(\frac{\pi s}{2}\right)
\Gamma(1-s)
\zeta(1-s).
\label{eq:functional_equation_zeta}
\end{equation}

The symmetry encoded by
\eqref{eq:functional_equation}
plays a central role in the study of the zeros of
$\zeta(s)$.

\subsection{Special Values and Basic Identities}
\label{subsec:special_values}

Among the classical identities are

\begin{equation}
\zeta(2)
=
\frac{\pi^{2}}{6},
\label{eq:zeta2}
\end{equation}

and more generally

\begin{equation}
\zeta(2m)
=
(-1)^{m+1}
\frac{B_{2m}(2\pi)^{2m}}
{2(2m)!},
\qquad
m\ge 1,
\label{eq:zeta_even_values}
\end{equation}

where
$B_{2m}$
denotes the Bernoulli numbers.

These identities are standard and may be found in
\cite{Titchmarsh1986}
or
\cite{Edwards2001}.

\subsection{The Zeros of the Zeta Function}
\label{subsec:zeta_zeros}

The zeros of
$\zeta(s)$
are traditionally divided into two classes.

\begin{definition}
\label{def:trivial_zeros}

The trivial zeros are located at

\begin{equation}
s=-2,-4,-6,\ldots.
\label{eq:trivial_zeros}
\end{equation}

\end{definition}

All remaining zeros are called nontrivial zeros.

\begin{definition}
\label{def:nontrivial_zero}

A nontrivial zero is a complex number

\begin{equation}
\rho
=
\beta+i\gamma
\label{eq:rho_definition}
\end{equation}

such that

\begin{equation}
\zeta(\rho)=0,
\label{eq:rho_zero_condition}
\end{equation}

with

\begin{equation}
0<\beta<1.
\label{eq:critical_strip}
\end{equation}

\end{definition}

The region
\eqref{eq:critical_strip}
is known as the critical strip.

The functional equation implies that the nontrivial zeros occur in
symmetric sets

\begin{equation}
\rho,
\quad
\overline{\rho},
\quad
1-\rho,
\quad
1-\overline{\rho}.
\label{eq:zero_symmetry}
\end{equation}

\begin{definition}
\label{def:zeta_zero_counting_function}

For \(T>0\), let

\begin{equation}
N_\zeta(T)
=
\#\{\rho=\beta+i\gamma:\zeta(\rho)=0,\ 0<\beta<1,\ 0<\gamma\leq T\},
\label{eq:zeta_zero_counting_function}
\end{equation}

where zeros are counted with multiplicity.

\end{definition}

The Riemann--von Mangoldt formula gives the asymptotic estimate

\begin{equation}
N_\zeta(T)
=
\frac{T}{2\pi}
\log\!\left(
\frac{T}{2\pi}
\right)
-
\frac{T}{2\pi}
+
O(\log T).
\label{eq:riemann_von_mangoldt}
\end{equation}

See
\cite{Titchmarsh1986}
or
\cite{Edwards2001}.

\section{The Riemann Hypothesis}
\label{sec:riemann_hypothesis}

This section contains standard material.

\begin{conjecture}[Riemann Hypothesis]
\label{conj:RH}

Every nontrivial zero of the Riemann zeta function satisfies

\begin{equation}
\Re(\rho)
=
\frac12.
\label{eq:RH}
\end{equation}

\end{conjecture}

Assuming
\eqref{eq:RH},
all nontrivial zeros can be written as

\begin{equation}
\rho_n
=
\frac12
+
i\gamma_n.
\label{eq:RH_zeros}
\end{equation}

The sequence of positive ordinates is written as

\begin{equation}
0<\gamma_1
\leq
\gamma_2
\leq
\gamma_3
\leq
\cdots,
\label{eq:gamma_ordering}
\end{equation}

with multiplicities counted. This sequence is the principal object
studied throughout this work.

The Riemann Hypothesis is one of the most important open problems in
mathematics and has deep implications for the distribution of prime
numbers.
Comprehensive discussions may be found in
\cite{Titchmarsh1986},
\cite{Edwards2001},
and
\cite{IwaniecKowalski2004}.

\section{Uniform Distribution of Zeta Ordinates}
\label{sec:uniform_distribution_zeta_ordinates}

The study of exponential sums over the ordinates of the nontrivial
zeros is closely related to the theory of uniform distribution modulo
one.

The fundamental tool is Weyl's criterion.

\begin{theorem}[Weyl Criterion]
\label{thm:weyl_criterion}

Let $(x_n)$ be a sequence of real numbers.

Then $(x_n)$ is uniformly distributed modulo one if and only if

\begin{equation}
\lim_{N\to\infty}
\frac{1}{N}
\sum_{n=1}^{N}
e^{2\pi i m x_n}
=
0
\qquad
(m\neq 0).
\label{eq:weyl_criterion}
\end{equation}

\end{theorem}

A natural question is therefore whether the sequence

\begin{equation}
\alpha \gamma_n
\label{eq:alpha_gamma_sequence}
\end{equation}

is uniformly distributed modulo one.

The first systematic investigation of this problem appears in the work
of Rademacher \cite{Rademacher1956}.

Rademacher studied the exponential sums

\begin{equation}
A_\alpha(T)
=
\sum_{0<\gamma\le T}
e^{2\pi i \alpha \gamma},
\label{eq:rademacher_sum}
\end{equation}

and asked whether cancellation occurs for general values of
$\alpha$.

Subsequent work by Hlawka established a remarkably strong
equidistribution result.

\begin{theorem}[Hlawka, 1975]
\label{thm:hlawka}

For every nonzero real number $\alpha$,
the sequence

\begin{equation}
\alpha \gamma_n
\label{eq:hlawka_sequence}
\end{equation}

is uniformly distributed modulo one.

Equivalently,

\begin{equation}
\lim_{N\to\infty}
\frac{1}{N}
\sum_{n=1}^{N}
e^{2\pi i m \alpha \gamma_n}
=
0
\qquad
(m\neq 0).
\label{eq:hlawka_weyl}
\end{equation}

\end{theorem}

A particularly important choice for the present work is

\begin{equation}
\alpha
=
\frac{1}{2\pi}.
\label{eq:alpha_two_pi_inverse}
\end{equation}

In that case,

\begin{equation}
x_n
=
\frac{\gamma_n}{2\pi}
\label{eq:gamma_over_two_pi}
\end{equation}

is uniformly distributed modulo one.

Applying Weyl's criterion with harmonic \(m=1\) gives

\begin{align}
\lim_{N\to\infty}
\frac{1}{N}
\sum_{n=1}^{N}
e^{2\pi i x_n}
&=
0
\label{eq:weyl_application_first_harmonic}
\\
&=
\lim_{N\to\infty}
\frac{1}{N}
\sum_{n=1}^{N}
e^{i\gamma_n}.
\label{eq:weyl_application_SN}
\end{align}

Therefore

\begin{equation}
\lim_{N\to\infty}
\frac{1}{N}
\sum_{n=1}^{N}
e^{i\gamma_n}
=
0.
\label{eq:SN_normalized_cancellation}
\end{equation}

In terms of the exponential sum

\begin{equation}
S_N
=
\sum_{n=1}^{N}
e^{i\gamma_n},
\label{eq:SN_definition_hlawka}
\end{equation}

this becomes

\begin{equation}
\frac{S_N}{N}
\longrightarrow
0.
\label{eq:SN_over_N_limit}
\end{equation}

Thus the classical theory predicts normalized cancellation of the
sequence \(S_N\).

It is important to emphasize that
\eqref{eq:SN_over_N_limit}
does not imply

\begin{equation}
S_N
\longrightarrow
0.
\label{eq:SN_not_implied}
\end{equation}

The asymptotic statement concerns only the average value of
\(S_N\) relative to \(N\).

The possibility that the trajectory \(S_N\) may exhibit additional
geometric structure is one of the motivations for the investigations
carried out in this paper.

The exceptional frequencies that survive averaging are described by the
Landau--Gonek theory discussed in Subsection~\ref{sec:prime_power_frequencies}.

\subsection{The Landau--Gonek Formula}
\label{sec:prime_power_frequencies}

Unlike the metric theory of Hlawka,
the Landau--Gonek formula studies special frequencies associated with
prime powers.

For a fixed integer \(x>1\), define the pure phase sum

\begin{equation}
S_x(T)
=
\sum_{0<\gamma\le T}
x^{\,i\gamma}.
\label{eq:SxT_definition}
\end{equation}

The behaviour of \eqref{eq:SxT_definition} depends strongly on the
arithmetic nature of \(x\).

\begin{definition}[von Mangoldt Function]
\label{def:von_mangoldt_function}

The von Mangoldt function is defined by

\begin{equation}
\Lambda(n)
=
\begin{cases}
\log p,
&
n=p^{m},
\quad p \text{ prime},
\quad m\ge 1,
\\[1ex]
0,
&
\text{otherwise}.
\end{cases}
\label{eq:von_mangoldt}
\end{equation}

\end{definition}

The function \(\Lambda(n)\) therefore detects prime powers and assigns
to each prime power \(p^m\) the weight \(\log p\).

It plays a central role in the explicit formulae relating prime numbers
and nontrivial zeros of the Riemann zeta function
\cite{Landau1911,Titchmarsh1986,IwaniecKowalski2004}.

The classical Landau formula is naturally stated for the zero sum

\begin{equation}
L_x(T)
=
\sum_{0<\gamma\le T}
x^\rho.
\label{eq:landau_zero_sum}
\end{equation}

For a fixed integer \(x>1\), Landau's formula gives

\begin{equation}
L_x(T)
=
-
\frac{T}{2\pi}
\Lambda(x)
+
O(\log T).
\label{eq:landau_zero_sum_asymptotic}
\end{equation}

In particular, when \(x\) is a prime power,

\begin{equation}
x=p^m,
\label{eq:prime_power_condition}
\end{equation}

the von Mangoldt term in \eqref{eq:landau_zero_sum_asymptotic} is
nonzero.

Assuming the Riemann Hypothesis, each nontrivial zero has the form

\begin{equation}
\rho
=
\frac{1}{2}
+
i\gamma.
\label{eq:RH_zero_decomposition}
\end{equation}

Therefore

\begin{equation}
x^\rho
=
\sqrt{x}\,x^{i\gamma}.
\label{eq:xrho_phase_decomposition}
\end{equation}

Combining \eqref{eq:landau_zero_sum_asymptotic} and
\eqref{eq:xrho_phase_decomposition}, for a fixed integer \(x>1\),
the pure phase sum satisfies, under the Riemann Hypothesis,

\begin{equation}
S_x(T)
=
-
\frac{T}{2\pi}
\frac{\Lambda(x)}{\sqrt{x}}
+
O(\log T).
\label{eq:landau_phase_sum_asymptotic}
\end{equation}


The modern quantitative form of this phenomenon is due to Gonek
\cite{Gonek1984},
who refined Landau's original argument using the explicit formula.

\begin{theorem}[Landau--Gonek, phase form under RH]
\label{thm:landau_gonek_formula}

For frequencies associated with prime powers,
the pure phase sum \eqref{eq:SxT_definition} admits, under the Riemann
Hypothesis, a main term proportional to the von Mangoldt function
divided by \(\sqrt{x}\).

In the fixed-\(x\) form,

\begin{equation}
\sum_{0<\gamma\le T}
x^{\,i\gamma}
=
-
\frac{T}{2\pi}
\frac{\Lambda(x)}{\sqrt{x}}
+
O(\log T).
\label{eq:landau_gonek_asymptotic}
\end{equation}

\end{theorem}

Consequently,
prime powers define distinguished resonant frequencies for exponential
sums over the ordinates of the nontrivial zeros.

This phenomenon stands in sharp contrast with the metric
equidistribution results discussed in the previous section.

\begin{remark}
\label{rem:prime_power_resonances}

Equation
\eqref{eq:landau_phase_sum_asymptotic}
\cite{Landau1911,Gonek1984}
shows that exponential sums over the ordinates of the zeta zeros possess
distinguished resonant frequencies indexed by prime powers.

The factor \(x^{-1/2}\) appears because
\eqref{eq:SxT_definition} is a pure phase sum in \(x^{i\gamma}\),
whereas Landau's original formula is naturally stated for the zero sum
in \(x^\rho\).

From the viewpoint of Weyl equidistribution,
these frequencies are exceptional:
they generate a nonzero main term at the Landau--Gonek scale,
while remaining compatible with normalized equidistribution.

The contrast between the generic behaviour described by
Theorem~\ref{thm:hlawka}
and the resonant behaviour described by
Theorem~\ref{thm:landau_gonek_formula}
is one of the central themes underlying the study of exponential sums
over zeta ordinates.

\end{remark}


\subsection{Equidistribution Modulo One}
\label{sec:equidistribution_modulo_one}

The ordinates \(\gamma_n\) of the nontrivial zeros do not form a
periodic or arithmetic sequence.

Given a fixed nonzero real parameter \(\alpha\), the equidistribution
problem concerns the sequence of fractional parts

\begin{equation}
\{\alpha\gamma_n\}.
\label{eq:scaled_ordinate_fractional_parts}
\end{equation}

Equidistribution modulo one means that, for every interval
\(I\subset[0,1]\),

\begin{equation}
\frac{1}{N}
\#\{1\le n\le N:\{\alpha\gamma_n\}\in I\}
\longrightarrow
|I|
\qquad
(N\to\infty).
\label{eq:equidistribution_interval_definition}
\end{equation}

The distribution modulo one of sequences associated with the ordinates
of the nontrivial zeros was studied by Rademacher and Hlawka
\cite{Rademacher1956,Hlawka1975}.
Later, Fujii obtained quantitative results on the uniformity of this
distribution
\cite{Fujii1978}.

Equivalently, after projection to the unit circle, the phases

\begin{equation}
e^{2\pi i\alpha\gamma_n}
\label{eq:scaled_ordinate_phase}
\end{equation}

have no asymptotically preferred arc.

The particular case

\begin{equation}
\alpha
=
\frac{1}{2\pi}
\label{eq:two_pi_inverse_scaling}
\end{equation}

corresponds to reduction of the ordinates modulo \(2\pi\), since

\begin{equation}
e^{2\pi i\alpha\gamma_n}
=
e^{i\gamma_n}.
\label{eq:two_pi_scaling_phase_identity}
\end{equation}

Thus, equidistribution of

\begin{equation}
\frac{\gamma_n}{2\pi}
\label{eq:ordinate_mod_two_pi_sequence}
\end{equation}

modulo one is equivalent to equidistribution of the phases
\(e^{i\gamma_n}\) on the unit circle.

Equidistribution is therefore a global angular distribution statement.
It does not yet address the local spacing, pair correlations, or
spectral interpretations of the zero ordinates.


\section{The Exponential Sum $S_N$}
\label{sec:SN_definition}

This section introduces the main objects investigated in this paper.

\subsection{Data Source}
\label{subsec:data_source}

The numerical computations reported in this work use the first

\begin{equation}
N
=
2\,001\,052
\label{eq:number_of_zeros}
\end{equation}

nontrivial zeros of the Riemann zeta function.

The ordinates were obtained from the publicly available tables of
Andrew Odlyzko:

\begin{equation}
\texttt{https://www-users.cse.umn.edu/\~{}odlyzko/zeta\_tables/zeros6}.
\label{eq:odlyzko_source}
\end{equation}

The dataset provides high-precision values of the ordinates

\begin{equation}
\gamma_n,
\qquad
n=1,\ldots,N.
\label{eq:gamma_dataset}
\end{equation}

\subsection{Definition of the Phase Walk}
\label{subsec:phase_walk}

For each ordinate
\(\gamma_n\)
we associate the unit complex phase

\begin{equation}
z_n
=
e^{i\gamma_n}.
\label{eq:phase_definition}
\end{equation}

\begin{definition}
\label{def:SN}

The cumulative exponential sum is defined by

\begin{equation}
S_N
=
\sum_{k=1}^{N}
e^{i\gamma_k}.
\label{eq:SN_definition}
\end{equation}

\end{definition}

Equivalently,

\begin{equation}
S_N
=
\sum_{k=1}^{N}
\cos(\gamma_k)
+
i
\sum_{k=1}^{N}
\sin(\gamma_k).
\label{eq:SN_real_imag}
\end{equation}

The sequence

\begin{equation}
S_1,S_2,\ldots,S_N
\label{eq:SN_trajectory}
\end{equation}

defines a trajectory in the complex plane obtained by successively
adding the unit vectors
\(e^{i\gamma_k}\).

\begin{definition}
\label{def:CN}

The centroid of the phase-walk trajectory is defined by

\begin{equation}
C_N
=
\frac{1}{N}
\sum_{n=1}^{N}
S_n.
\label{eq:CN_definition}
\end{equation}

\end{definition}

Using the definition of \(S_n\), the centroid can also be written as

\begin{equation}
C_N
=
\frac{1}{N}
\sum_{n=1}^{N}
\sum_{k=1}^{n}
e^{i\gamma_k}.
\label{eq:CN_double_sum}
\end{equation}

Interchanging the order of summation gives the weighted form

\begin{equation}
C_N
=
\frac{1}{N}
\sum_{k=1}^{N}
(N-k+1)e^{i\gamma_k}.
\label{eq:CN_weighted}
\end{equation}

Thus \(C_N\) is not the simple average of the phases.
It is the average of the partial sums \(S_n\), or equivalently a
triangularly weighted average of the increments \(e^{i\gamma_k}\).

The geometric and distributional properties of the trajectory
\((S_n)\), together with the behaviour of its centroid \(C_N\),
constitute the central subject of the present paper.

\begin{remark}
\label{rem:equidistribution_not_centroid}

Classical results of Weyl, Hlawka and subsequent authors imply that the
normalized exponential sums satisfy

\begin{equation}
\frac{1}{N}
\sum_{n=1}^{N}
e^{i\gamma_n}
\longrightarrow
0.
\label{eq:normalized_phase_cancellation}
\end{equation}

Equivalently,

\begin{equation}
S_N=o(N).
\label{eq:SN_sublinear_growth}
\end{equation}

Thus, from the point of view of first-order averaging,
the phases \(e^{i\gamma_n}\) exhibit asymptotic cancellation.

However, this fact alone does not imply that

\begin{equation}
S_N
\longrightarrow
0,
\label{eq:SN_zero_limit}
\end{equation}

nor does it imply that the centroid \(C_N\) is located at the origin.

Indeed, by \eqref{eq:CN_definition},
the centroid is an average of the partial sums \(S_n\), not a simple
average of the phases themselves.

Equivalently, by \eqref{eq:CN_weighted},
the centroid is a triangularly weighted average of the increments
\(e^{i\gamma_k}\).

Consequently, the existence of asymptotic phase cancellation is fully
compatible with the presence of nontrivial geometric structure in the
trajectory \((S_n)\), including the possibility of a nonzero limiting
centroid.

The purpose of the present investigation is to characterize the
geometry of the trajectory \((S_n)\), including its centroid,
phase distribution, and large-scale structure, and to explore how
these features relate to the classical theory of exponential sums over
the ordinates of the nontrivial zeros.

\end{remark}

\begin{remark}
\label{rem:original_object}

The definition of the phase walk
\eqref{eq:SN_definition}
is the starting point for the original numerical investigations
presented in this work.

\end{remark}

\section{Equivalent Product Forms for \(S_N\) and \(C_N\)}
\label{sec:product_forms_SN_CN}

This section gives exact algebraic reformulations of the cumulative
sum \(S_N\) and of the centroid \(C_N\) in terms of finite products
associated with the phases \(e^{i\gamma_k}\).

\subsection{A Finite Product Identity}
\label{subsec:finite_product_identity}

We first record the elementary identity on which the product
formulation is based.

\begin{lemma}[Logarithmic coefficients of a finite product]
\label{lem:finite_product_logarithmic_coefficients}

Let \(a_1,\ldots,a_N\) be complex numbers, and define

\begin{equation}
P(z)
=
\prod_{k=1}^{N}
(1-a_k z).
\label{eq:finite_product_definition}
\end{equation}

For \(|z|\) sufficiently small,

\begin{equation}
\log P(z)
=
-
\sum_{m\ge 1}
\frac{z^m}{m}
\sum_{k=1}^{N}
a_k^m.
\label{eq:finite_product_logarithmic_expansion}
\end{equation}

In particular,

\begin{equation}
-
\left.
\frac{d}{dz}
\log P(z)
\right|_{z=0}
=
\sum_{k=1}^{N}
a_k.
\label{eq:first_logarithmic_coefficient}
\end{equation}

\end{lemma}

\begin{proof}

Since \(P(z)\) is a finite product,

\begin{equation}
\log P(z)
=
\sum_{k=1}^{N}
\log(1-a_k z).
\label{eq:finite_product_logarithm_sum}
\end{equation}

Using

\begin{equation}
\log(1-w)
=
-
\sum_{m\ge 1}
\frac{w^m}{m}
\label{eq:log_one_minus_w_expansion}
\end{equation}

for \(|w|\) sufficiently small, we obtain

\begin{align}
\log P(z)
&=
-
\sum_{k=1}^{N}
\sum_{m\ge 1}
\frac{(a_k z)^m}{m}
\label{eq:finite_product_expansion_first}
\\
&=
-
\sum_{m\ge 1}
\frac{z^m}{m}
\sum_{k=1}^{N}
a_k^m.
\label{eq:finite_product_expansion_second}
\end{align}

Differentiating at \(z=0\) gives
\eqref{eq:first_logarithmic_coefficient}.

\end{proof}
\subsection{The Phase Product and the Sum \(S_N\)}
\label{subsec:phase_product_SN}

We now apply Lemma~\ref{lem:finite_product_logarithmic_coefficients}
to the phases

\begin{equation}
a_k
=
e^{i\gamma_k}.
\label{eq:phase_product_ak}
\end{equation}

\begin{definition}[Phase product]
\label{def:phase_product}

The phase product associated with the first \(N\) ordinates is

\begin{equation}
P_N(z)
=
\prod_{k=1}^{N}
\left(1-z e^{i\gamma_k}\right).
\label{eq:phase_product_definition}
\end{equation}

\end{definition}

The zeros of \(P_N\) are

\begin{equation}
z
=
e^{-i\gamma_k},
\qquad
k=1,\ldots,N.
\label{eq:phase_product_zeros}
\end{equation}

Thus all zeros of \(P_N\) lie on the unit circle.

For each integer \(m\ge 1\), define the \(m\)-th phase sum by

\begin{equation}
S_N^{(m)}
=
\sum_{k=1}^{N}
e^{im\gamma_k}.
\label{eq:higher_phase_sum_definition}
\end{equation}

Then Lemma~\ref{lem:finite_product_logarithmic_coefficients} gives

\begin{equation}
\log P_N(z)
=
-
\sum_{m\ge 1}
\frac{S_N^{(m)}}{m}
z^m.
\label{eq:phase_product_logarithmic_expansion}
\end{equation}

In particular,

\begin{equation}
S_N
=
S_N^{(1)}.
\label{eq:SN_first_phase_mode}
\end{equation}

Equivalently,

\begin{equation}
S_N
=
-
\left.
\frac{d}{dz}
\log P_N(z)
\right|_{z=0}.
\label{eq:SN_logarithmic_derivative}
\end{equation}

Thus \(S_N\) is the first logarithmic coefficient of the phase product
\(P_N\).

\subsection{Polynomial Form and Elementary Symmetric Coefficients}
\label{subsec:elementary_symmetric_coefficients}

It is also useful to consider the monic polynomial

\begin{equation}
Q_N(z)
=
\prod_{k=1}^{N}
\left(z-e^{i\gamma_k}\right).
\label{eq:monic_phase_polynomial}
\end{equation}

Expanding \(Q_N\), one obtains

\begin{equation}
Q_N(z)
=
z^N
-
e_1^{(N)}z^{N-1}
+
e_2^{(N)}z^{N-2}
-
\cdots
+
(-1)^N e_N^{(N)}.
\label{eq:monic_phase_polynomial_expansion}
\end{equation}

Here

\begin{equation}
e_r^{(N)}
=
\sum_{1\le k_1<\cdots<k_r\le N}
e^{i(\gamma_{k_1}+\cdots+\gamma_{k_r})}.
\label{eq:elementary_symmetric_phase_coefficients}
\end{equation}

In particular,

\begin{equation}
e_1^{(N)}
=
\sum_{k=1}^{N}
e^{i\gamma_k}
=
S_N.
\label{eq:first_elementary_coefficient_SN}
\end{equation}

Hence

\begin{equation}
Q_N(z)
=
z^N
-
S_N z^{N-1}
+
e_2^{(N)}z^{N-2}
-
\cdots
+
(-1)^N e_N^{(N)}.
\label{eq:monic_phase_polynomial_SN_coefficient}
\end{equation}

This identifies \(S_N\) as the first elementary symmetric coefficient
of the phases \(e^{i\gamma_k}\).

\subsection{Product Form of the Centroid}
\label{subsec:product_form_centroid}

The centroid was defined by

\begin{equation}
C_N
=
\frac{1}{N}
\sum_{n=1}^{N}
S_n.
\label{eq:CN_definition_product_section}
\end{equation}

For each \(n\le N\), define the partial phase product

\begin{equation}
P_n(z)
=
\prod_{k=1}^{n}
\left(1-z e^{i\gamma_k}\right).
\label{eq:partial_phase_product_definition}
\end{equation}

By \eqref{eq:SN_logarithmic_derivative}, applied with \(N\) replaced by
\(n\), one has

\begin{equation}
S_n
=
-
\left.
\frac{d}{dz}
\log P_n(z)
\right|_{z=0}.
\label{eq:Sn_partial_product_logarithmic_derivative}
\end{equation}

Substitution into \eqref{eq:CN_definition_product_section} gives

\begin{equation}
C_N
=
-
\frac{1}{N}
\sum_{n=1}^{N}
\left.
\frac{d}{dz}
\log P_n(z)
\right|_{z=0}.
\label{eq:CN_sum_logarithmic_derivatives}
\end{equation}

By linearity of differentiation,

\begin{equation}
C_N
=
-
\left.
\frac{d}{dz}
\left(
\frac{1}{N}
\sum_{n=1}^{N}
\log P_n(z)
\right)
\right|_{z=0}.
\label{eq:CN_average_logarithmic_derivative}
\end{equation}
\begin{definition}[Accumulated phase product]
\label{def:accumulated_phase_product}

The accumulated phase product is defined by

\begin{equation}
\mathcal{P}_N(z)
=
\prod_{n=1}^{N}
P_n(z).
\label{eq:accumulated_phase_product_definition}
\end{equation}

\end{definition}

Since

\begin{equation}
\log \mathcal{P}_N(z)
=
\sum_{n=1}^{N}
\log P_n(z),
\label{eq:accumulated_phase_product_logarithm}
\end{equation}

we obtain the exact product identity

\begin{equation}
C_N
=
-
\frac{1}{N}
\left.
\frac{d}{dz}
\log \mathcal{P}_N(z)
\right|_{z=0}.
\label{eq:CN_accumulated_product_logarithmic_derivative}
\end{equation}

Thus \(C_N\) is the normalized first logarithmic coefficient of the
accumulated phase product \(\mathcal{P}_N\).

\subsection{The Triangular Product}
\label{subsec:triangular_phase_product}

Using the definition of \(P_n\),

\begin{equation}
\mathcal{P}_N(z)
=
\prod_{n=1}^{N}
\prod_{k=1}^{n}
\left(1-z e^{i\gamma_k}\right).
\label{eq:accumulated_phase_product_double_product}
\end{equation}

Changing the order of multiplication gives

\begin{equation}
\mathcal{P}_N(z)
=
\prod_{k=1}^{N}
\left(1-z e^{i\gamma_k}\right)^{N-k+1}.
\label{eq:triangular_phase_product}
\end{equation}

Combining
\eqref{eq:CN_accumulated_product_logarithmic_derivative}
and
\eqref{eq:triangular_phase_product},
one obtains

\begin{equation}
C_N
=
-
\frac{1}{N}
\left.
\frac{d}{dz}
\log
\prod_{k=1}^{N}
\left(1-z e^{i\gamma_k}\right)^{N-k+1}
\right|_{z=0}.
\label{eq:CN_triangular_product_logarithmic_derivative}
\end{equation}

Differentiating explicitly gives

\begin{equation}
C_N
=
\frac{1}{N}
\sum_{k=1}^{N}
(N-k+1)e^{i\gamma_k}.
\label{eq:CN_triangular_weighted_sum}
\end{equation}

This is the triangular weighted form of the centroid.

\subsection{Triangular Harmonic Interpretation}
\label{subsec:triangular_harmonic_interpretation}

The usual empirical measure of the phases is

\begin{equation}
\nu_N
=
\frac{1}{N}
\sum_{k=1}^{N}
\delta_{e^{i\gamma_k}}.
\label{eq:uniform_empirical_phase_measure}
\end{equation}

Its first moment is

\begin{equation}
\int_{\mathbb T}
z\,d\nu_N(z)
=
\frac{1}{N}
\sum_{k=1}^{N}
e^{i\gamma_k}
=
\frac{S_N}{N}.
\label{eq:uniform_empirical_phase_first_moment}
\end{equation}

The centroid \(C_N\), by contrast, corresponds to the triangularly
weighted finite measure

\begin{equation}
\mu_N^{\triangle}
=
\frac{1}{N}
\sum_{k=1}^{N}
(N-k+1)
\delta_{e^{i\gamma_k}}.
\label{eq:triangular_phase_measure}
\end{equation}

Its first moment is

\begin{equation}
\int_{\mathbb T}
z\,d\mu_N^{\triangle}(z)
=
\frac{1}{N}
\sum_{k=1}^{N}
(N-k+1)e^{i\gamma_k}
=
C_N.
\label{eq:triangular_phase_measure_first_moment}
\end{equation}

Therefore \(C_N\) is the first harmonic of a triangularly weighted phase
measure, whereas \(S_N/N\) is the first harmonic of the uniformly
weighted empirical phase measure.

\subsection{Local Analytic Interpretation at the Origin}
\label{subsec:local_analytic_interpretation_origin}

From \eqref{eq:phase_product_logarithmic_expansion},

\begin{equation}
\log P_N(z)
=
-
S_N z
+
O(z^2)
\qquad
(z\to 0).
\label{eq:phase_product_local_expansion}
\end{equation}

Thus

\begin{equation}
S_N
=
-
\left.
\frac{d}{dz}
\log P_N(z)
\right|_{z=0}.
\label{eq:SN_local_analytic_interpretation}
\end{equation}

Similarly, from
\eqref{eq:CN_accumulated_product_logarithmic_derivative},

\begin{equation}
C_N
=
-
\frac{1}{N}
\left.
\frac{d}{dz}
\log \mathcal{P}_N(z)
\right|_{z=0}.
\label{eq:CN_local_analytic_interpretation}
\end{equation}

Consequently, \(S_N\) and \(C_N\) are local logarithmic derivatives at
the origin of two finite products associated with the phases
\(e^{i\gamma_k}\): the ordinary phase product \(P_N\) and the accumulated
phase product \(\mathcal{P}_N\), respectively.

\begin{remark}
\label{rem:product_form_summary}

The identities above are exact finite-\(N\) identities.

They do not depend on the Riemann Hypothesis, on random matrix models,
or on probabilistic assumptions about the zeros.

They show that \(S_N\) is the first logarithmic coefficient of the
ordinary phase product \(P_N\), while \(C_N\) is the normalized first
logarithmic coefficient of the accumulated product
\(\mathcal{P}_N\).

\end{remark}


\section{The Non-Resonant Triangular Landau--Gonek Structure of \(S_N\) and \(C_N\)}
\label{sec:non_resonant_triangular_lg_structure}

This section records rigorous structural identities for the ordinal sum
\(S_N\), the height-truncated sums \(S_X(T)\), and the centroid \(C_N\).
The Landau--Gonek input is literature-based; the terminology describing
\(C_N\) as a triangular non-resonant transform is an original structural
interpretation of the identities below.

\subsection{Ordinal and height truncations}
\label{subsec:ordinal_and_height_truncations}

The sum \(S_N\) is an ordinal truncation: \(N\) is a positive integer
and counts the first \(N\) ordinates in the ordered list
\eqref{eq:gamma_ordering}. By contrast, height truncations are cut off by
an analytic height \(T\) on the ordinate axis. The counting function
\(N_\zeta(T)\), defined in \eqref{eq:zeta_zero_counting_function},
translates the height \(T\) into the corresponding ordinal number of
zeros. Thus \(N\) is not identified with \(T\).

\begin{definition}
\label{def:continuous_height_phase_sum}

For \(X\in\mathbb{R}_{>0}\) and \(T>0\), define

\begin{equation}
S_X(T)
=
\sum_{0<\gamma\leq T}X^{i\gamma},
\label{eq:continuous_height_phase_sum}
\end{equation}

where zeros are counted with multiplicity and

\begin{equation}
X^{i\gamma}:=\exp(i\gamma\log X)
\label{eq:continuous_power_real_log}
\end{equation}

with the real logarithm of \(X\).

\end{definition}

\begin{theorem}[Height-to-ordinal bridge]
\label{thm:height_to_ordinal_bridge}

For every \(T>0\),

\begin{equation}
S_e(T)
=
\sum_{0<\gamma\leq T}e^{i\gamma}
=
\sum_{k=1}^{N_\zeta(T)}e^{i\gamma_k}
=
S_{N_\zeta(T)}.
\label{eq:height_to_ordinal_bridge}
\end{equation}

\end{theorem}

\begin{proof}
The ordered sequence \eqref{eq:gamma_ordering} lists the positive
ordinates of the nontrivial zeros with multiplicity. By definition,
\(N_\zeta(T)\) is exactly the number of entries in this list whose
ordinate is at most \(T\). Therefore the height-truncated multiset
\(\{\gamma:0<\gamma\leq T\}\), counted with multiplicity, is the same
multiset as \(\{\gamma_k:1\leq k\leq N_\zeta(T)\}\). Summing
\(e^{i\gamma}\) over these two identical multisets gives the first two
equalities in \eqref{eq:height_to_ordinal_bridge}. The last equality is
precisely the ordinal definition \eqref{eq:SN_definition} with
\(N=N_\zeta(T)\).
\end{proof}

\begin{corollary}[Ordinal recovery away from ambiguity]
\label{cor:ordinal_recovery_from_height}

For every height \(T>0\) such that \(N_\zeta(T)=N\),

\begin{equation}
S_N=S_e(T).
\label{eq:ordinal_recovery_from_height}
\end{equation}

\end{corollary}

\begin{proof}
Substitute the hypothesis \(N_\zeta(T)=N\) into
\eqref{eq:height_to_ordinal_bridge}.
\end{proof}

\subsection{Continuous phase sums}
\label{subsec:continuous_phase_sums}

The parameter \(X\) in \eqref{eq:continuous_height_phase_sum} is a
positive real continuous frequency. The definition
\eqref{eq:continuous_power_real_log} avoids any multivalued complex
logarithm.

\begin{theorem}[The continuous frequency \(X=e\)]
\label{thm:e_continuous_frequency_phase}

At the continuous frequency \(X=e\), the height-truncated phase sum is

\begin{equation}
S_e(T)
=
\sum_{0<\gamma\leq T}e^{i\gamma}.
\label{eq:e_continuous_frequency_phase_sum}
\end{equation}

\end{theorem}

\begin{proof}
Since \(\log e=1\), definition \eqref{eq:continuous_power_real_log}
gives

\begin{equation}
e^{i\gamma}:=\exp(i\gamma\log e)=\exp(i\gamma).
\label{eq:e_phase_identity}
\end{equation}

Inserting this identity into \eqref{eq:continuous_height_phase_sum} with
\(X=e\) proves \eqref{eq:e_continuous_frequency_phase_sum}.
\end{proof}

\begin{remark}
\label{rem:continuous_not_arithmetic_frequency}

The frequency \(X=e\) belongs to \(\mathbb{R}_{>0}\). It is used here as
a continuous phase frequency, not as an arithmetic argument of the von
Mangoldt function.
\end{remark}

\subsection{Arithmetic Landau--Gonek resonances}
\label{subsec:arithmetic_landau_gonek_resonances}

Landau--Gonek resonances are arithmetic. In this subsection
\(X\in\mathbb{R}_{>0}\) denotes a continuous frequency, while
\(x\in\mathbb{N}\), \(x>1\), denotes an arithmetic frequency.

\begin{definition}
\label{def:prime_power_resonance_set}

The arithmetic resonance set is

\begin{equation}
\mathcal R
=
\{p^m:\ p\ \text{prime},\ m\in\mathbb N\}.
\label{eq:prime_power_resonance_set}
\end{equation}

\end{definition}

For fixed \(x\in\mathbb N\), \(x>1\), the Landau--Gonek formula under RH
has the phase form

\begin{equation}
\sum_{0<\gamma\leq T}x^{i\gamma}
=
-
\frac{T}{2\pi}
\frac{\Lambda(x)}{\sqrt{x}}
+
O(\log T).
\label{eq:landau_gonek_new_section_phase_form}
\end{equation}

The main term is supported on the prime-power set \(\mathcal R\).

\begin{theorem}[Non-resonance of \(e\)]
\label{thm:e_non_resonant}

The continuous frequency \(X=e\) is not an arithmetic Landau--Gonek
resonance:

\begin{equation}
e\notin\mathcal R.
\label{eq:e_not_in_resonance_set}
\end{equation}

\end{theorem}

\begin{proof}
Every element of \(\mathcal R\) is a positive integer, so
\(\mathcal R\subset\mathbb N\). The number \(e\) is irrational, hence
\(e\notin\mathbb N\). Therefore \(e\notin\mathcal R\).
\end{proof}

\begin{remark}
\label{rem:no_mangoldt_at_e}

The statement is not that the von Mangoldt function is evaluated at
\(e\). The correct formulation is that \(X=e\) is a positive real
continuous frequency, but it is not an arithmetic frequency in the
prime-power resonance set \(\mathcal R\).
\end{remark}

\subsection{The centroid as a triangular non-resonant transform}
\label{subsec:centroid_triangular_non_resonant_transform}

The centroid identity is purely ordinal and contains no height variable.
It follows from the definitions in Section~\ref{sec:SN_definition}.

\begin{theorem}[Triangular centroid identity]
\label{thm:triangular_centroid_identity}

For every positive integer \(N\),

\begin{equation}
C_N
=
\frac{1}{N}\sum_{n=1}^{N}S_n
=
\frac{1}{N}\sum_{n=1}^{N}\sum_{k=1}^{n}e^{i\gamma_k}
=
\frac{1}{N}\sum_{k=1}^{N}(N-k+1)e^{i\gamma_k}.
\label{eq:triangular_centroid_identity}
\end{equation}

\end{theorem}

\begin{proof}
The first equality is the definition \eqref{eq:CN_definition}. Replacing
\(S_n\) by its defining sum gives the second equality, which is
\eqref{eq:CN_double_sum}. In the double sum, the term
\(e^{i\gamma_k}\) appears exactly for those indices \(n\) with
\(k\leq n\leq N\). There are \(N-k+1\) such indices. Collecting equal
terms gives the final expression, already recorded in
\eqref{eq:CN_weighted}.
\end{proof}

\begin{corollary}
\label{cor:centroid_triangular_non_resonant_phase_sum}

\(C_N\) is a triangularly weighted non-resonant phase sum associated with
the continuous frequency \(X=e\).

\end{corollary}

\begin{proof}
By Theorem~\ref{thm:triangular_centroid_identity}, \(C_N\) is a
triangularly weighted sum of the phases \(e^{i\gamma_k}\). By
Theorem~\ref{thm:e_continuous_frequency_phase}, these phases are the
phases generated by the continuous frequency \(X=e\). By
Theorem~\ref{thm:e_non_resonant}, this frequency is not in the
arithmetic resonance set \(\mathcal R\). Hence the triangularly weighted
ordinal sum defining \(C_N\) is non-resonant in the Landau--Gonek
prime-power sense.
\end{proof}

\subsection{Equidistribution consequences}
\label{subsec:equidistribution_consequences_non_resonant}

Hlawka's theorem, applied with \(\alpha=1/(2\pi)\), implies the
normalized cancellation already stated in \eqref{eq:SN_over_N_limit}.
The next consequence records the same fact in the notation of this
section.

\begin{corollary}[Non-resonant normalized cancellation]
\label{cor:non_resonant_normalized_cancellation}

The ordinal phase sums at the continuous frequency \(X=e\) satisfy

\begin{equation}
\frac{S_N}{N}\longrightarrow 0
\qquad
(N\to\infty).
\label{eq:non_resonant_normalized_cancellation}
\end{equation}

\end{corollary}

\begin{proof}
Equation \eqref{eq:SN_over_N_limit} is obtained from Hlawka's
equidistribution theorem and Weyl's criterion with
\(\alpha=1/(2\pi)\). Since \(e^{i\gamma_k}\) is the phase associated with
\(X=e\) by Theorem~\ref{thm:e_continuous_frequency_phase}, the result is
exactly \eqref{eq:non_resonant_normalized_cancellation}.
\end{proof}

\begin{remark}
\label{rem:centroid_not_forced_by_equidistribution}

The cancellation in \eqref{eq:non_resonant_normalized_cancellation}
concerns the uniformly weighted average of the increments
\(e^{i\gamma_k}\). It does not force the triangularly weighted centroid
\(C_N\) to vanish, because the weights in
\eqref{eq:triangular_centroid_identity} depend on the ordinal position
\(k\).
\end{remark}

\subsection{Structural conclusion}
\label{subsec:structural_conclusion_non_resonant}

The preceding results separate three layers of structure. First,
\(S_N\) is an ordinal sum, while \(S_X(T)\) is a height-truncated sum.
Second, \(N_\zeta(T)\) is the bridge from analytic height to ordinal
length. Third, the special phase \(e^{i\gamma}\) is generated by the
continuous frequency \(X=e\), and this frequency is not a prime-power
Landau--Gonek resonance.

The paper identifies \(C_N\) as a triangular non-resonant Landau--Gonek
type transform at the continuous frequency \(X=e\).

\begin{thebibliography}{99}

\bibitem{Riemann1859}
B. Riemann,
\emph{\"Uber die Anzahl der Primzahlen unter einer gegebenen Gr\"osse},
Monatsberichte der Berliner Akademie,
1859,
671--680.

Reprinted in:

B. Riemann,
\emph{Gesammelte Mathematische Werke},
Dover Publications,
1953.

\bibitem{Titchmarsh1986}
E. C. Titchmarsh,
\emph{The Theory of the Riemann Zeta-Function},
Second Edition,
revised by D. R. Heath-Brown,
Oxford University Press,
1986.

\bibitem{Edwards2001}
H. M. Edwards,
\emph{Riemann's Zeta Function},
Dover Publications,
2001.

\bibitem{IwaniecKowalski2004}
H. Iwaniec and E. Kowalski,
\emph{Analytic Number Theory},
American Mathematical Society,
2004.

\bibitem{Weyl1916}
H. Weyl,
\emph{\"Uber die Gleichverteilung von Zahlen mod. Eins},
Mathematische Annalen
\textbf{77} (1916),
313--352.
doi:10.1007/BF01475864.

\bibitem{Rademacher1956}
H. Rademacher,
\emph{On the distribution of the values of the zeta function},
Colloquium Mathematicum
\textbf{4} (1956),
79--81.

\bibitem{Hlawka1975}
E. Hlawka,
\emph{\"Uber die Gleichverteilung gewisser Folgen,
die mit den Nullstellen der Zetafunktion zusammenh\"angen},
Sitzungsberichte der
\"Osterreichischen Akademie der Wissenschaften,
Mathematisch-Naturwissenschaftliche Klasse,
Abteilung II
\textbf{184} (1975),
459--471.

\bibitem{Landau1911}
E. Landau,
\emph{\"Uber die Nullstellen der Zetafunktion},
Mathematische Annalen
\textbf{71} (1911),
548--564.
doi:10.1007/BF01456881.

\bibitem{Gonek1984}
S. M. Gonek,
\emph{A Formula of Landau and Its Applications to the Zeros of the Zeta Function},
in:
\emph{A Tribute to Emil Grosswald:
Number Theory and Related Analysis},
Contemporary Mathematics
\textbf{143},
American Mathematical Society,
1993,
395--413.

\bibitem{Borwein2008}
P. Borwein,
S. Choi,
B. Rooney,
A. Weirathmueller,
\emph{The Riemann Hypothesis:
A Resource for the Afficionado and Virtuoso Alike},
Springer,
2008.

\bibitem{Odlyzko1987}
A. M. Odlyzko,
\emph{On the Distribution of Spacings Between Zeros of the Zeta Function},
Mathematics of Computation
\textbf{48} (1987),
273--308.
doi:10.2307/2007886.

\bibitem{OdlyzkoTables}
A. M. Odlyzko,
\emph{Tables of Zeros of the Riemann Zeta Function},
available at
\url{https://www-users.cse.umn.edu/~odlyzko/zeta_tables/}.

\bibitem{FordZaharescu2005}
K. Ford and A. Zaharescu,
\emph{Unnormalized Differences Between Zeros of the Riemann Zeta Function},
Acta Arithmetica
\textbf{117} (2005),
163--176.
doi:10.4064/aa117-2-5.

\bibitem{FordSoundararajanZaharescu2024}
K. Ford,
K. Soundararajan,
A. Zaharescu,
\emph{The Distribution of the Ordinates of the Zeros of the Riemann Zeta Function},
arXiv:2407.14173,
2024.

\bibitem{Fujii1978}
A. Fujii,
\emph{On the uniformity of the distribution of the zeros of the
Riemann zeta function},
Journal f\"ur die reine und angewandte Mathematik
\textbf{302} (1978),
167--205.

\end{thebibliography}

\end{document}

